% Chapter 2: 分布回归的极小极大速率
\chapter{分布回归的极小极大速率}\label{ch:chapter2}

% ----------------------------------------------------------
\section{问题背景与研究动机}\label{sec:ch2-background}

\textbf{分布回归}（distribution regression）的目标是：
给定协变量 $X \in \mathbb{R}^{D_X}$，估计响应变量 $Y \in \mathbb{R}^{D_Y}$
的条件分布 $\mu^*_{Y|X=x}$。
与传统的标量回归或分类不同，分布回归恢复的是完整的条件分布，
从而可以刻画条件变异性、偏度、多峰性和不确定性等丰富信息，
在生物医学、气候建模和计量经济学中有广泛应用。

传统非参数密度回归方法（如核密度估计）存在两个根本性局限：

\begin{enumerate}
    \item \textbf{存在性假设}：要求 $\mu^*_{Y|X=x}$ 关于 $\mathbb{R}^{D_Y}$ 上
          Lebesgue 测度绝对连续——当 $Y$ 含离散分量或位于高维空间中的低维流形上时，
          该假设失效。
    \item \textbf{维度灾难}：环境维度 $D_X, D_Y$ 主导收敛速率，
          对现代高维数据（图像、文本、点云）效率极低。
\end{enumerate}

本文~\citep{tang2025regression} 放弃了 $\mu^*_{Y|X=x}$ 须有密度的假设，
转而假设 $\mu^*_{Y|X=x}$ 支撑于某个低维流形 $\mathcal{M}_{Y|x}$ 上，
从而能够同时处理奇异分布和流形结构数据。

\textbf{与条件生成模型的关系}：
条件生成对抗网络（cGAN）、条件扩散模型、条件归一化流
均可视为隐式分布回归方法——它们不显式估计密度，
而是通过采样逼近真实条件分布。
本文建立的极小极大理论为这些方法的理论性能提供了基准。

% ----------------------------------------------------------
\section{数学框架}\label{sec:ch2-framework}

\subsection{基本设定}\label{sec:ch2-setup}

设
\begin{itemize}
    \item 协变量 $X \sim \mu^*_X$，支撑于 $\mathbb{R}^{D_X}$ 中的 $d_X$ 维子流形
          $\mathcal{M}_X$；
    \item 给定 $X=x$，响应 $Y \sim \mu^*_{Y|X=x}$，支撑于 $\mathbb{R}^{D_Y}$ 中的
          $d_Y$ 维 $\beta_Y$-光滑子流形 $\mathcal{M}_{Y|x}$；
    \item 数据 $\{(X_i, Y_i)\}_{i=1}^n$ 为来自 $\mu^* = \mu^*_X \mu^*_{Y|X}$
          的 $n$ 个 i.i.d. 样本。
\end{itemize}

当 $\mathcal{M}_{Y|x}$ 依赖于 $x$ 时，假设族
$\{\mathcal{M}_{Y|x} : x \in \mathcal{M}_X\}$ 是
$(\beta_Y, \beta_X)$-光滑的（见\cref{def:covdep-manifold}）。
条件密度 $u^*(y|x)$ 关于 $\mathcal{M}_{Y|x}$ 的体积测度是
$\alpha_Y$-光滑（关于 $y$）和 $\alpha_X$-光滑（关于 $x$）的。

\textbf{协变量空间维数假设}（比流形假设更弱）：
\begin{Definition}[协变量空间的 Minkowski 维数]\label{def:minkowski-dim}
    称拓扑空间 $\mathcal{M}_X \subset \mathbb{R}^{D_X}$ 的
    Minkowski 维数至多为 $d_X$（记 $\mathcal{M}_X \in \mathcal{M}_X(D_X, d_X, L)$），
    若 $\mathcal{M}_X \subset \mathbb{B}_{\mathbb{R}^{D_X}}(\mathbf{0}, L)$，
    且对任意 $0 < \varepsilon \leq 1$，
    $\mathcal{M}_X$ 的最大 $\varepsilon$-packing 基数至多为 $L\varepsilon^{-d_X}$。
\end{Definition}

该定义不要求 $\mathcal{M}_X$ 光滑或有正 reach，比子流形假设更宽松。

% ----------------------------------------------------------
\subsection{$\gamma$-Hölder IPM 度量}\label{sec:ch2-ipm}

本文采用积分概率度量（IPM，亦称对抗损失）来量化估计误差。
给定光滑度参数 $\gamma \geq 0$，定义 $\gamma$-\textbf{Hölder IPM} 为
\citep[亦见][]{tang2023twosample}：
\begin{equation}\label{eq:gamma-holder-ipm}
    d_\gamma(\mu, \nu) \;=\; \sup_{f \in \mathcal{H}_1^\gamma(\mathbb{R}^{D_Y})}
    \left|\int f(y)\,\mathrm{d}\mu - \int f(y)\,\mathrm{d}\nu\right|,
\end{equation}
其中 $\mathcal{H}_1^\gamma(\mathbb{R}^{D_Y})$ 为单位半径的 $\gamma$-Hölder 函数类。

特别地：
\begin{itemize}
    \item $\gamma = 0$：$d_0$ 退化为\textbf{全变差距离} $d_{\mathrm{TV}}$
         （取常值 1 的函数即可）；
    \item $\gamma = 1$：$d_1$ 等价于 \textbf{1-Wasserstein 距离} $W_1$
          （由 Kantorovich-Rubinstein 定理）。
\end{itemize}

参数 $\gamma$ 调控度量的"敏感度"：
$\gamma$ 越大，测试函数越光滑，越能忽略局部细节（如支撑不对齐），
而对全局概率质量差异更敏感；$\gamma$ 越小，$d_\gamma$ 对支撑错位和密度局部起伏越敏感。

评估条件分布估计量 $\widehat{\mu}_{Y|X}$ 的整体质量采用
\textbf{期望 $\gamma$-Hölder IPM}：
\begin{equation}\label{eq:expected-ipm}
    \mathbb{E}_{\mu^*_X}\!\left[d_\gamma\!\big(\mu^*_{Y|X},\; \widehat{\mu}_{Y|X}\big)\right]
    \;=\; \int_{\mathcal{M}_X} d_\gamma\!\big(\mu^*_{Y|X=x},\; \widehat{\mu}_{Y|X=x}\big)\,
          \mathrm{d}\mu^*_X(x).
\end{equation}

\begin{Remark}\label{rem:gamma-meaning}
    $\gamma$ 的选择对应于应用中不同的关注点：
    \begin{itemize}
        \item $\gamma = 0$（全变差）：关注整体概率质量分配，
              对支撑几何不敏感；
        \item $\gamma = 1$（Wasserstein）：关注概率质量在空间中的
              运输成本，对支撑几何敏感；
        \item $\gamma \in (0,1)$：介于二者之间，
              同时平衡支撑恢复与密度估计的权衡。
    \end{itemize}
    \citep[亦见][]{tang2025regression} 的理论覆盖所有 $\gamma \geq 0$，
    从而统一了全变差距离和 Wasserstein 距离的分析。
\end{Remark}

% ----------------------------------------------------------
\section{Regime 1：经典密度回归}\label{sec:ch2-regime1}

本节考虑最简设定：响应空间无流形结构，
即 $\mathcal{M}_{Y|x} = \mathbb{R}^{D_Y}$，
$\mu^*_{Y|x}$ 对 Lebesgue 测度绝对连续。
此为经典密度回归设定。

\subsection{Regime 1 的分布族}\label{sec:ch2-regime1-def}

\begin{Definition}[Regime 1：欧氏响应空间]\label{def:regime1}
    给定维度 $D_Y, D_X \in \mathbb{N}_+$、$d_X \in \mathbb{N} \cap [0, D_X]$、
    光滑度参数 $\alpha_Y, \alpha_X \in (0, \infty)$ 和常数 $L > 0$，
    定义分布族
    \[
        \mathcal{P}_1^* = \mathcal{P}_1^*(D_Y, D_X, d_X, \alpha_Y, \alpha_X, L)
    \]
    为满足以下条件的联合分布 $\mu^* = \mu^*_X\mu^*_{Y|X}$ 构成的集合：
    \begin{enumerate}
        \item $\mu^*_X$ 的支撑 $\mathcal{M}_X \in \mathcal{M}_X(D_X, d_X, L)$，
              $\mu^*_Y$ 的支撑含于 $\mathbb{B}_{\mathbb{R}^{D_Y}}(\mathbf{0}, L)$；
        \item 对任意 $x \in \mathcal{M}_X$，$\mu^*_{Y|x}$ 有密度函数
              $u^*(\cdot|x)$，且
              $u^*(y|x) \in \overline{\mathcal{H}}^{\alpha_Y, \alpha_X}_L
                 (\mathbb{R}^{D_Y}, \mathcal{M}_X)$
              （即 $u^*(\cdot|y)$ 对 $y$ 为 $\alpha_Y$-光滑，
              对 $x$ 为 $\alpha_X$-光滑；见\cref{def:marginal-smooth}）。
    \end{enumerate}
\end{Definition}

\begin{Remark}\label{rem:dX-zero}
    $d_X = 0$ 对应 $X$ 为离散变量（取有限个值），
    此时 $\alpha_X$ 的光滑性条件自动满足，
    极小极大速率简化为 $n^{-(\alpha_Y+\gamma)/(2\alpha_Y+D_Y)}$，
    即~\citep{Li2022} 的无条件密度估计结果。
\end{Remark}

\subsection{极小极大速率}\label{sec:ch2-regime1-rate}

\begin{Theorem}[Regime 1 的极小极大速率~\citep{tang2025regression}]\label{def:thm-regime1}
    对任意 $\gamma \geq 0$，存在常数 $L_0$，
    使得当 $L, n \geq L_0$ 时，
    \begin{equation}\label{eq:regime1-minimax}
        \begin{aligned}
            & C\!\left(
                n^{-\frac{\alpha_X}{2\alpha_X + d_X}}
                + n^{-\frac{\alpha_Y + \gamma}{2\alpha_Y + D_Y + \frac{\alpha_Y}{\alpha_X}d_X}}
            \right) \\
            \leq\;&
            \inf_{\widehat{\mu}_{Y|X}}\;
            \sup_{\mu^* \in \mathcal{P}_1^*}
            \;
            \mathbb{E}_{\mu^{*,\otimes n}}\!\left[
                \mathbb{E}_{\mu^*_X}\!\left[
                    d_\gamma\!\big(\mu^*_{Y|X},\; \widehat{\mu}_{Y|X}\big)
                \right]
            \right] \\
            \leq\;&
            C_1\;\widetilde{\mathcal{O}}\!\left(
                n^{-\frac{\alpha_X}{2\alpha_X + d_X}}
                + n^{-\frac{\alpha_Y + \gamma}{2\alpha_Y + D_Y + \frac{\alpha_Y}{\alpha_X}d_X}}
            \right),
        \end{aligned}
    \end{equation}
    其中 $C, C_1$ 为与 $n$ 无关的常数，$\widetilde{\mathcal{O}}$ 忽略对数因子。
\end{Theorem}

\textbf{速率的两项意义}：
\begin{enumerate}
    \item \textbf{第一项} $n^{-\alpha_X/(2\alpha_X+d_X)}$：
          对应 $\alpha_X$-光滑条件均值函数在 $L_2$ 损失下的
          经典极小极大速率~\citep{Stone1982}。
          $\gamma$ 越大，该项越可能主导——
          因为光滑测试函数抹平了条件密度的局部起伏，
          $d_\gamma$ 主要响应 $X$ 与 $Y$ 之间的全局依赖趋势。
    \item \textbf{第二项} $n^{-(\alpha_Y+\gamma)/(2\alpha_Y+D_Y+\alpha_Y d_X/\alpha_X)}$：
          反映非参数条件密度估计的固有难度。
          $\alpha_Y$ 越大（密度越光滑）或 $d_X, D_Y$ 越小，速率越快；
          $\gamma$ 越大，$d_\gamma$ 对局部细节越不敏感，
          密度估计难度降低。
\end{enumerate}

\textbf{与已有文献的关系}：
\citep{tang2025regression} 的结果改进了~\citep{Bilodeau2023} 的上界
（在 $\gamma=0$ 时给出更紧的全变差速率）
并拓展了~\citep{Li2022} 的结果（仅覆盖 $\alpha_X \in [0,1]$）；
~\citep{tang2024diffusion} 中条件扩散模型的收敛速率
在 $\gamma=1$ 时与本定理上界匹配，
结合下界可知条件扩散模型在期望 1-Wasserstein 距离下是极小极大最优的。

% ----------------------------------------------------------
\section{Regime 2：响应空间为协变量无关流形}\label{sec:ch2-regime2}

本节将设定推广到 $Y$ 支撑于低维流形上的情形，
但假设 $\mathcal{M}_{Y|x} = \mathcal{M}_Y$ 不依赖于 $x$——
此时问题从密度回归变为\textbf{分布回归}。

\subsection{Regime 2 的分布族}\label{sec:ch2-regime2-def}

流形上的密度函数须关于流形体积测度而非 Lebesgue 测度定义，
且流形 $\mathcal{M}_Y$ 本身未知。

\begin{Definition}[Regime 2：协变量无关的流形响应空间]\label{def:regime2}
    给定维度参数 $D_Y, d_Y, D_X \in \mathbb{N}_+$、$d_X \in \mathbb{N}_0$、
    光滑度 $\beta_Y, \alpha_Y, \alpha_X > 0$、
    函数 $g: \mathbb{R}_+ \to \mathbb{R}_+$ 和常数 $\tau, \tau_1, L > 0$，
    定义分布族
    \[
        \mathcal{P}_2^* = \mathcal{P}_2^*(D_Y, D_X, d_Y, d_X,
            \beta_Y, \alpha_Y, \alpha_X, \tau, \tau_1, g, L)
    \]
    为满足以下条件的 $\mu^* = \mu^*_X\mu^*_{Y|X}$ 构成的集合：
    \begin{enumerate}
        \item $\mathcal{M}_X \in \mathcal{M}_X(D_X, d_X, L)$；
        \item 对任意 $x \in \mathcal{M}_X$，$\mu^*_{Y|x}$ 支撑于
              与 $x$ 无关的子流形 $\mathcal{M}_Y$ 上，
              且有相对于 $\mathcal{M}_Y$ 体积测度的密度函数
              $u^*(\cdot|x)$，
              其中 $\mathcal{M}_Y \in \mathcal{M}_{\tau,\tau_1,L}^{\beta_Y}(d_Y, D_Y)$
             （$d_Y$ 维、$\beta_Y$-光滑子流形），
              $u^* \in \overline{\mathcal{H}}_L^{\alpha_Y, \alpha_X}
                 (\mathcal{M}_Y, \mathcal{M}_X)$；
        \item 下界条件：
              对任意 $x_0 \in \mathcal{M}_X$、$y_0 \in \mathcal{M}_Y$ 和 $0 < r \leq 1$，
              $\mu^*_X(\mathbb{B}_{\mathcal{M}_X}(x_0, r)) \geq g(r)r^{d_X}/L$
              且 $\mu^*_{Y|x_0}(\mathbb{B}_{\mathcal{M}_Y}(y_0, r))
                 \geq g(r)r^{d_Y}/L$。
    \end{enumerate}
\end{Definition}

条件 $\beta_Y \geq 2$ 保证 $\mathcal{M}_Y$ 曲率有界；
$\beta_Y \geq \alpha_Y + 1$ 保证 $\alpha_Y$ 光滑性与局部坐标卡的选择无关
（流形内在的光滑性定义）。

\subsection{极小极大速率}\label{sec:ch2-regime2-rate}

\begin{Theorem}[Regime 2 的极小极大速率~\citep{tang2025regression}]\label{def:thm-regime2}
    对任意 $\gamma > 0$ 和 $\beta_Y \geq 2 \vee (\alpha_Y + 1)$，
    存在常数 $L_0$，使得当 $L, \tau, \tau_1, n \geq L_0$ 时，
    \begin{equation}\label{eq:regime2-minimax}
        \begin{aligned}
            & C\!\left(
                n^{-\frac{\alpha_X}{2\alpha_X + d_X}}
                + n^{-\frac{\alpha_Y + \gamma}{2\alpha_Y + d_Y + \frac{\alpha_Y}{\alpha_X}d_X}}
                + n^{-\frac{\gamma\beta_Y}{d_Y}}
            \right) \\
            \leq\;&
            \inf_{\widehat{\mu}_{Y|X}}\;
            \sup_{\mu^* \in \mathcal{P}_2^*}
            \;
            \mathbb{E}_{\mu^{*,\otimes n}}\!\left[
                \mathbb{E}_{\mu^*_X}\!\left[
                    d_\gamma\!\big(\mu^*_{Y|X},\; \widehat{\mu}_{Y|X}\big)
                \right]
            \right] \\
            \leq\;&
            C_1\;(\log n)^3\;
            \widetilde{\mathcal{O}}\!\left(
                n^{-\frac{\alpha_X}{2\alpha_X + d_X}}
                + n^{-\frac{\alpha_Y + \gamma}{2\alpha_Y + d_Y + \frac{\alpha_Y}{\alpha_X}d_X}}
                + n^{-\frac{\gamma\beta_Y}{d_Y}}
            \right).
        \end{aligned}
    \end{equation}
\end{Theorem}

\textbf{与 Regime 1 的对比}：
Regime 2 的速率多了第三项 $n^{-\gamma\beta_Y/d_Y}$，
反映了对未知流形 $\mathcal{M}_Y$ 进行支撑估计的固有难度。
当 $\gamma$ 较小（$d_\gamma$ 对支撑几何敏感）时，
该项主导误差——这正是 $d_\gamma$ 度量奇异分布的独特之处。
当 $d_X = 0$（$X$ 离散）时，前两项简化为
$n^{-1/2} + n^{-(\alpha_Y+\gamma)/(2\alpha_Y+d_Y)}$，
与~\citep{tang2023twosample} 的无条件流形分布估计结果一致。

\textbf{关键现象——$\gamma$ 决定三项竞争}：
由~\cref{fig:regime2-phase}（取自原文 Figure 1），
随着 $\gamma$ 从 0 增大，
Regime 2 的主导误差项经历两次相变：

\begin{itemize}
    \item \textbf{第一次相变}（$\gamma$ 较小）：
          $d_\gamma$ 对支撑错位敏感，
          主导项为流形估计项 $n^{-\gamma\beta_Y/d_Y}$
          $\Longleftrightarrow$ 条件密度估计项
          $n^{-(\alpha_Y+\gamma)/(2\alpha_Y+d_Y+\alpha_Y d_X/\alpha_X)}$；
    \item \textbf{第二次相变}（$\gamma$ 较大）：
          $d_\gamma$ 进一步光滑化，
          主导项变为 $n^{-\alpha_X/(2\alpha_X+d_X)}$——
          此时问题退化为 $\alpha_X$-光滑条件均值估计。
\end{itemize}

\begin{figure}[htbp]
    \centering
    % Figures are in the PDF directory
    \caption{Regime 2 中极小极大速率的相图（取自~\citep{tang2025regression}）：
            (a) 固定 $\alpha_X$、$d_X$、$d_Y$、$\beta_Y$，极小极大速率主导项
            随 $\alpha_Y$ 和 $\gamma$ 的变化；
            (b) 固定 $d_X$、$d_Y$、$\beta_Y$，主导项随 $\alpha_X$ 和 $\gamma$ 的变化。
            两个临界 $\gamma$ 值决定三项之间的竞争。}\label{fig:regime2-phase}
\end{figure}

\begin{Remark}\label{rem:TV-fail}
    在 Regime 2 中，由于 $\widehat{\mu}_{Y|X}$ 与 $\mu^*_{Y|x}$
    几乎必然相互奇异（支撑在不同流形上），
    全变差距离恒为 1，无法反映估计质量——
    因此 Regime 2 须要求 $\gamma > 0$（即使用非平凡光滑测试函数）。
    这凸显了 IPM 框架比全变差距离更适用于流形支撑数据的重要性。
\end{Remark}

% ----------------------------------------------------------
\section{Regime 3：协变量依赖流形}\label{sec:ch2-regime3}

本节研究最一般设定：响应流形 $\mathcal{M}_{Y|x}$ 随 $x$ 光滑变化——
这将问题从单一流形估计推广为\textbf{流形回归}。

\subsection{协变量依赖流形的形式定义}\label{sec:ch2-covdep-manifold}

\begin{Definition}[$(\beta_Y, \beta_X)$-光滑子流形族~\citep{tang2025regression}]\label{def:covdep-manifold}
    称 $\{\mathcal{M}_{Y|x} : x \in \mathcal{M}_X\}$ 为
    $(\beta_Y, \beta_X)$-光滑的 $d_Y$ 维子流形族，
    若存在常数 $\tau, \tau_1, L > 0$ 使得：
    对每个 $x \in \mathcal{M}_X$，$\mathcal{M}_{Y|x} \in
       \mathcal{M}_{\tau,\tau_1,L}^{\beta_Y}(d_Y, D_Y)$，
    且对任意 $x, x' \in \mathcal{M}_X$，
    \[
        \mathbb{H}(\mathcal{M}_{Y|x},\; \mathcal{M}_{Y|x'})
        \;\leq\; L\,\|x - x'\|^{\beta_X/\beta_Y},
    \]
    其中 $\mathbb{H}$ 为 Hausdorff 距离。
    即族的光滑性由 $\beta_X$ 控制——
    $\beta_X$ 越大，$\mathcal{M}_{Y|x}$ 随 $x$ 变化越慢，
    信息跨不同 $x$ 值的共享越容易。
\end{Definition}

\textbf{直观例子}：
考虑人脸图像数据 $Y$ 以年龄和性别 $X$ 为条件的情形。
对固定的 $X$，$Y$ 的支撑自然位于某低维图像子流形；
而不同年龄段对应的流形差异较大，
这正构成协变量依赖流形族。

% ----------------------------------------------------------
\subsection{流形回归的极小极大速率}\label{sec:ch2-manifold-regression}

为建立 Regime 3b 的下界，先研究纯流形回归问题——
仅估计 $\{\mathcal{M}_{Y|x}\}$ 而不计密度。

\begin{Theorem}[流形回归的极小极大速率~\citep{tang2025regression}]\label{def:thm-manifold-regression}
    假设 $\beta_Y \geq 2$ 和 $\beta_Y \geq \beta_X$，
    则存在常数 $L_0$，使得当 $L, \tau, \tau_1, n \geq L_0$ 时，
    \begin{equation}\label{eq:manifold-regression-rate}
        C\,n^{-\frac{1}{\frac{d_Y}{\beta_Y} + \frac{d_X}{\beta_X}}}
        \;\leq\;
        \inf_{\{\widehat{\mathcal{M}}_{Y|x}\}}
        \sup_{\mu^* \in \mathcal{P}^*}
        \;
        \mathbb{E}_{\mu^{*,\otimes n}}\!\left[
            \sup_{x \in \mathcal{M}_X}
            \mathbb{H}\!\big(\mathcal{M}_{Y|x},\; \widehat{\mathcal{M}}_{Y|x}\big)
        \right]
        \;\leq\;
        C_1\!\left(\frac{n}{\log n}\right)^{\!
            -\frac{1}{\frac{d_Y}{\beta_Y} + \frac{d_X}{\beta_X}}}.
    \end{equation}
\end{Theorem}

与单一流形估计的极小极小速率
$n^{-1/(d_Y/\beta_Y)}$~\citep{AamariLevrard2019} 相比，
流形回归多了 $d_X/\beta_X$ 项，反映了估计整个流形族的额外复杂度。
$\beta_X$ 越大（$\mathcal{M}_{Y|x}$ 随 $x$ 变化越慢），
分母越大，速率越快——这与直觉一致。

\textbf{估计策略}：推广了~\citep{AamariLevrard2019} 的局部多项式估计量，
对每个数据点 $w_k = (X_k, Y_k)$，
选取满足 $\|Y - Y_k\| \leq h_1$ 和 $\|X - X_k\| \leq h_2$
的邻域样本，其中
$h_1 \asymp (\log n/n)^{\frac{\beta_X}{d_Y\beta_X + d_X\beta_Y}}$、
$h_2 \asymp (\log n/n)^{\frac{\beta_Y}{d_Y\beta_X + d_X\beta_Y}}$，
然后通过最小化重建损失学习局部多项式估计量。

% ----------------------------------------------------------
\subsection{Regime 3b：协变量依赖流形上的分布回归}\label{sec:ch2-regime3b}

现在回到完整的分布回归问题，但此时 $\mathcal{M}_{Y|x}$ 依赖 $x$。

\begin{Theorem}[Regime 3b 的极小极大速率~\citep{tang2025regression}]\label{def:thm-regime3b}
    对任意 $\gamma > 0$，
    若 $\beta_Y \geq 2 \vee (\alpha_Y + 1) \vee \beta_X$、
    $\beta_X \geq \alpha_X + \frac{\alpha_X}{\alpha_Y}$ 且 $\alpha_Y \geq \alpha_X$，
    则存在常数 $L_0$，使得当 $L, \tau, \tau_1, n \geq L_0$ 时，
    \begin{equation}\label{eq:regime3b-minimax}
        \begin{aligned}
            & C\!\left(
                n^{-\frac{\alpha_X}{2\alpha_X + d_X}}
                + n^{-\frac{\alpha_Y + \gamma}{2\alpha_Y + d_Y + \frac{\alpha_Y}{\alpha_X}d_X}}
                + n^{-\frac{\gamma}{\frac{d_Y}{\beta_Y} + \frac{d_X}{\beta_X}}}
            \right) \\
            \leq\;&
            \inf_{\widehat{\mu}_{Y|X}}\;
            \sup_{\mu^* \in \mathcal{P}_3^*}
            \;
            \mathbb{E}_{\mu^{*,\otimes n}}\!\left[
                \mathbb{E}_{\mu^*_X}\!\left[
                    d_\gamma\!\big(\mu^*_{Y|X},\; \widehat{\mu}_{Y|X}\big)
                \right]
            \right] \\
            \leq\;&
            C_1\;\widetilde{\mathcal{O}}\!\left(
                n^{-\frac{\alpha_X}{2\alpha_X + d_X}}
                + n^{-\frac{\alpha_Y + \gamma}{2\alpha_Y + d_Y + \frac{\alpha_Y}{\alpha_X}d_X}}
                + n^{-\frac{\gamma}{\frac{d_Y}{\beta_Y} + \frac{d_X}{\beta_X}}}
            \right).
        \end{aligned}
    \end{equation}
\end{Theorem}

\textbf{与 Regime 2 的关键区别}：
Regime 3b 的第三项为
$n^{-\gamma/(d_Y/\beta_Y + d_X/\beta_X)}$，
而非 Regime 2 中的 $n^{-\gamma\beta_Y/d_Y}$。
这一替换将单一流形估计的复杂度
$d_Y/\beta_Y$ 推广为流形回归的复杂度
$d_Y/\beta_Y + d_X/\beta_X$——
几何参数 $d_Y, \beta_Y, d_X, \beta_X$ 共同决定收敛速率。

\textbf{几何参数的物理意义}：
\begin{itemize}
    \item $d_Y$ 越小：响应流形本征维度越低，估计越容易；
    \item $\beta_Y$ 越大：$\mathcal{M}_{Y|x}$ 越光滑，
          其几何特征越规则，支撑估计越容易；
    \item $d_X$ 越小：协变量空间复杂度越低；
    \item $\beta_X$ 越大：$\mathcal{M}_{Y|x}$ 随 $x$ 变化越慢，
          允许跨 $x$ 值的信息共享更充分。
\end{itemize}

\begin{Remark}\label{rem:regime-comparison}
    三个 Regime 的极小极大速率可统一理解为三项的取大：
    \begin{enumerate}
        \item $n^{-\alpha_X/(2\alpha_X+d_X)}$：协变量效应；
        \item $n^{-(\alpha_Y+\gamma)/(2\alpha_Y+d_Y+\alpha_Y d_X/\alpha_X)}$：密度效应；
        \item 流形支撑效应（Regime 依赖）:
              Regime 2 为 $n^{-\gamma\beta_Y/d_Y}$，
              Regime 3b 为 $n^{-\gamma/(d_Y/\beta_Y+d_X/\beta_X)}$。
    \end{enumerate}
    当 $\gamma$ 足够大时，前两项之一主导；
    当 $\gamma$ 较小时，支撑估计项主导——
    这正是 $d_\gamma$ 度量奇异分布的核心挑战。
\end{Remark}

% ----------------------------------------------------------
\section{极小极大最优估计量}\label{sec:ch2-estimators}

本文不仅建立了极小极小下界，还提出了
融合对抗学习（adversarial learning）与同步最小二乘的
\textbf{混合估计量}，在各 Regime 下均达上界（至多对数因子）。

\subsection{Euclid 响应空间的估计量（Regime 1）}\label{sec:ch2-estimator-regime1}

估计量基于\textbf{小波展开}和\textbf{局部多项式回归}。
对条件密度 $u^*(y|x)$ 在 $y$ 方向做小波分解，
在 $x$ 方向做局部多项式逼近。

\textbf{核心思想}：
对每个小波基 $\psi \in \Psi_j^{D_Y}$ 在尺度 $j$，估计条件期望
$\mathbb{E}_{\mu^*_{Y|x}}[\psi(Y)]$——
这本身是一个非参数回归问题，
但 $u^*(\cdot|x)$ 的 $\alpha_X$-光滑性允许用局部多项式有效逼近。

构造条件密度估计量：
\begin{equation}\label{eq:density-estimator}
    \widehat{u}(y|x) \;=\;
    \sum_{j=0}^{J}\;\sum_{\psi \in \Psi_j^{D_Y}}
    \widehat{u}_\psi(x)\;\psi(y),
    \qquad
    x \in \mathcal{M}_X,
\end{equation}
其中 $\widehat{u}_\psi(x)$ 通过局部多项式回归获得。
截断水平 $J$ 的选择：
\begin{equation}\label{eq:J-choice}
    J \;\asymp\;
    \frac{1}{2\alpha_Y + D_Y + d_X\frac{\alpha_Y}{\alpha_X}}
    \log_2\!\left(\frac{n}{\log n}\right).
\end{equation}

\textbf{与扩散模型的深层联系}：
~\citep{tang2025regression} 指出，
多分辨率小波分析与分数阶前向-后向扩散模型之间存在深刻类比：
小波分解中不同尺度 $j$ 捕获从全局趋势到局部细节的层次结构，
正好对应扩散模型中时间变量 $t$ 从 $0$ 到 $T$ 的演化过程。
条件扩散模型的每步去噪可视为在某一分辨率水平上的条件均值估计。

\begin{Theorem}[Regime 1 估计量的收敛速率~\citep{tang2025regression}]\label{def:thm-estimator-regime1}
    在 $\mathcal{P}_1^*$ 假设下，
    对任意 $\gamma \geq 0$，
    \cref{eq:density-estimator} 定义的条件密度估计量满足
    \begin{equation}\label{eq:estimator-regime1-rate}
        \mathbb{E}_{\mu^*_X}\!\left[
            d_\gamma\!\big(\mu^*_{Y|X},\; \widehat{\mu}_{Y|X}\big)
        \right]
        \;\leq\;
        C_\gamma\;\widetilde{\mathcal{O}}\!\left(
            n^{-\frac{\alpha_X}{2\alpha_X + d_X}}
            + n^{-\frac{\alpha_Y + \gamma}{2\alpha_Y + D_Y + \frac{\alpha_Y}{\alpha_X}d_X}}
        \right),
    \end{equation}
    以至少 $1 - n^{-1}$ 的概率成立。
    速率与\cref{def:thm-regime1} 上界匹配，故为极小极大最优。
\end{Theorem}

% ----------------------------------------------------------
\subsection{流形响应空间的估计量（Regime 2 和 3b）}\label{sec:ch2-estimator-manifold}

当 $\mathcal{M}_{Y|x}$ 为流形时，条件密度关于 Lebesgue 测度不存在，
须将问题重新表述为对条件期望泛函
\[
    \mathcal{T}^*(f, x) \;=\; \mathbb{E}_{\mu^*_{Y|x}}[f(Y)],
    \qquad
    f \in \mathcal{H}_1^\gamma(\mathbb{R}^{D_Y})
\]
的同步估计。

\textbf{两步走策略}：

\textbf{步骤 1（粗糙尺度）}：将小波函数截断到 $J$ 层，
估计 $\mathbb{E}_{\mu^*_{Y|x}}[f_J(Y)]$，
其中 $f_J$ 为 $f$ 的粗糙（小波）逼近。
由于 $f_J$ 变化缓慢，可将 $\mu^*_{Y|x}$ 视为有密度，
使用局部多项式回归构造条件密度估计量
（类似于 Regime 1，但须在流形坐标系中处理）。

关键改进：对每个 $j$ 水平，
同时估计所有小波系数的条件期望，
构建\textbf{联合均值回归问题}：
\begin{equation}\label{eq:joint-mean-regression}
    \widehat{S}_j^\dagger
    \;=\;
    \arg\min_{S \in \mathcal{S}_j^\dagger}\;
    \frac{1}{n}\sum_{i=1}^{n}\;\sum_{\psi \in \Psi_j^{D_Y}}
    \Big(
        2^{\frac{j(d_Y - D_Y)}{2}}\psi(Y_i)
        - S(\psi, X_i)
    \Big)^2,
\end{equation}
其中 $\mathcal{S}_j^\dagger$ 为特定的非可分逼近族。
与分别独立估计每个小波系数不同，
联合估计利用了流形结构固有的稀疏性——
只有部分小波函数在 $\mathcal{M}_{Y|x}$ 上有非零条件均值。

\textbf{步骤 2（精细尺度）}：对 $f_J^\perp$，
采用\textbf{显式流形估计}步骤。
学习编码器 $Q: \mathbb{R}^{D_Y} \to \mathbb{R}^{d_Y}$
和条件解码器 $G: \mathbb{R}^{d_Y} \times \mathbb{R}^{D_X} \to \mathbb{R}^{D_Y}$，
使得 $y \approx G(Q(y), x)$ 在局部成立。
将数据映射到潜在空间 $\mathbb{R}^{d_Y}$，
在其中进行密度回归。

\begin{Remark}\label{rem:encoder-decoder}
    这种编码器-解码器框架正是潜扩散模型（latent diffusion models）
    \citep{Rombach2022}、变分自编码器（VAE）\citep{Kingma2013}
    和 Wasserstein 自编码器\citep{Tolstikhin2017} 的理论基础——
    本文从统计估计的角度严格证明了该框架的理论最优性。
\end{Remark}

\textbf{对抗学习的角色}：
给定步骤 1-2 中的条件期望估计量 $\widehat{\mathcal{I}}(f, x)$，
通过对抗训练构造条件分布估计量：
\[
    \widehat{\mu}_{Y|x}^\gamma
    \;=\;
    \arg\min_{\mu \in \mathcal{P}_Y^*}
    \;\sup_{f \in \mathcal{H}_1^\gamma(\mathbb{R}^{D_Y})}
    \left|
        \mathbb{E}_\mu[f(y)] - \widehat{\mathcal{I}}(f, x)
    \right|,
\]
其中 $\mathcal{P}_Y^*$ 为响应空间的候选分布族。
对多尺度 $\gamma$ 同时优化，可得对所有 $\gamma > 0$ 同时最优的估计量。

% ----------------------------------------------------------
\section{技术要点}\label{sec:ch2-technical}

\subsection{下界证明技术}\label{sec:ch2-lower-bound-techniques}

极小极小下界的证明综合运用了：

\begin{enumerate}
    \item \textbf{Fano 不等式}：将估计问题规约到多假设检验问题，
          通过构造难以区分的分布对建立信息论下界。
    \item \textbf{Varshamov-Gilbert 引理}：构造足够分散的假设序列，
          保证相邻假设间的 KL 散度可控。
    \item \textbf{经验过程理论}：chaining、peeling、localization 等技术
          控制函数类的复杂度。
    \item \textbf{Poincaré 不等式}：用于控制流形上分布的 Fisher 信息。
    \item \textbf{定理 8（泛化 Fano 方法）}：
          论文发展的技术工具，
          将多分辨率小波基与 Fano 方法结合，
          同时处理粗糙和精细尺度误差。
\end{enumerate}

下界的三项分别对应三类难度来源：
协变量维数（通过 Fano + 局部多项式构造）、
响应维数（通过小波基构造）和流形支撑（通过流形构造）。

\subsection{上界证明技术}\label{sec:ch2-upper-bound-techniques}

\begin{enumerate}
    \item \textbf{小波分解}：将测试函数分解为粗糙和精细分量，
          分别处理；
    \item \textbf{局部多项式回归}：
          利用 $\alpha_X$-光滑性控制粗糙尺度的逼近误差；
    \item \textbf{单位分解（Partition of Unity）}：
          在无全局参数化的流形上构造局部估计量，
          并光滑拼接——这是~\citep{tang2022minimax} 发展的核心技术；
    \item \textbf{Talagrand 集中不等式}：
          控制经验过程偏离均值的高概率界；
    \item \textbf{联合估计 vs 独立估计}：
          Regime 2 估计量采用联合均值回归而非分别估计各小波系数，
          关键在于利用了流形结构的稀疏性——
          独立估计忽略了不同小波系数之间的几何依赖。
\end{enumerate}

\begin{Lemma}[联合 vs 独立估计的最优性]\label{def:lem-joint-vs-separate}
    在 Regime 2 中，令 $\widehat{u}_j(x)$ 为独立估计得到的
    各尺度条件密度，
    $\widehat{u}_j^\dagger(x)$ 为联合估计得到的条件密度。
    则
    \[
        \mathbb{E}\!\left[
            \big|\widehat{u}_j^\dagger(x) - u^*(y|x)\big|^2
        \right]
        \;\leq\;
        \mathbb{E}\!\left[
            \big|\widehat{u}_j(x) - u^*(y|x)\big|^2
        \right]
        \quad \text{（至多差常数因子）},
    \]
    且在某些构造下严格更优。
    详见~\citep[Appendix B.2.3]{tang2025regression}。
\end{Lemma}

% ----------------------------------------------------------
\section{与系列工作的联系}\label{sec:ch2-connections}

本章结果在 Tang-Rong 系列中的地位：

\begin{enumerate}
    \item \textbf{与~\citep{tang2022minimax}（Annals 2022）的关系}：
          极小极大理论框架相同（$\gamma$-Hölder IPM、
          三项取大结构），但本文从无条件的流形分布估计
          推广到协变量条件设定。
    \item \textbf{与~\citep{tang2023lowdim}（ICML 2023）的关系}：
          MLDMAE 方法是 Regime 2 上界的具体算法实现，
          验证了单位分解 + 混合局部估计策略的实践有效性。
    \item \textbf{与~\citep{tang2024diffusion}（AISTATS 2024）的关系}：
          条件扩散模型在 Wasserstein 距离（$\gamma=1$）
          下与 Regime 1 上界匹配，证明了扩散模型的极小极小最优性。
    \item \textbf{与~\citep{tang2025conditional}（arXiv 2025）的关系}：
          条件扩散模型在流形设定下的收敛速率
          与 Regime 3b 的极小极小速率对照——
          实践工具能否在一般 $\gamma$ 下达到最优仍是开放问题。
\end{enumerate}

% ----------------------------------------------------------
\section{开放问题}\label{sec:ch2-open}

基于本章结果，以下问题值得进一步研究：

\begin{enumerate}
    \item \textbf{$\gamma \in (0,1)$ 时的可计算估计量}：
          混合估计量在 $\gamma \neq 0, 1$ 时涉及无穷维对抗优化，
          实践中如何有效实现？IPM 判别器的逼近误差如何量化？
    \item \textbf{流形维数自适应}：
          本文假设 $d_X, d_Y$ 已知；
          如何构造对这些参数同时自适应的估计量？
          Lepski 方法能否扩展到分布回归设定？
    \item \textbf{$\alpha_X$-自适应}：
          极小极大速率依赖于 $\alpha_X, \alpha_Y$ 等光滑度参数，
          构造对所有光滑度同时自适应的估计量是重要挑战。
    \item \textbf{噪声观测情形}：
          本文考虑无噪声观测（$X_i, Y_i$ 精确采样）；
          测量噪声或标签噪声下的极小极小速率是什么？
    \item \textbf{流形回归与分布回归的统一理论}：
          Regime 3a（纯流形回归）和 Regime 3b（带密度）
          是否可以统一在一个更简洁的框架下？
    \item \textbf{深度生成模型与 IPM 的结合}：
          当前理论主要围绕 Wasserstein 距离和全变差距离；
          对 MMD、Stein 差异等更一般的 IPM，
          深度生成模型能否达到相应的极小极小速率？
\end{enumerate}

% ----------------------------------------------------------
\section{小结}\label{sec:ch2-summary}

本章系统介绍了~\citep{tang2025regression} 关于分布回归极小极小速率的理论。

\textbf{核心贡献}：

\begin{enumerate}
    \item \textbf{统一框架}：$\gamma$-Hölder IPM 统一了全变差距离和 Wasserstein 距离，
          使得同一套理论可分析多种概率度量下的估计问题。
    \item \textbf{三类 Regime}：根据 $\mathcal{M}_{Y|x}$ 与 $x$ 的关系，
          建立了从经典密度回归到协变量依赖流形回归的层次化理论，
          各 Regime 的极小极小速率清晰刻画。
    \item \textbf{几何参数决定复杂度}：
          收敛速率由 $d_X, d_Y, \beta_X, \beta_Y, \alpha_X, \alpha_Y, \gamma$
          等参数的特定组合决定，
          揭示了低维流形结构如何缓解维度灾难。
    \item \textbf{最优估计量}：提出了融合对抗学习与同步最小二乘的混合估计量，
          在各 Regime 下均达极小极小上界（至多对数因子），
          为条件生成模型提供了理论基础。
\end{enumerate}

\textbf{关键洞见}：
流形支撑与密度的联合估计是分布回归的核心挑战；
$\gamma$ 的选择决定了"关注支撑"与"关注密度"的权衡；
联合估计（而非分别估计各小波系数）是流形设定下达到极小极小最优的关键。
